\chapter{System Architecture and Design}
\label{ch:architecture}

This chapter presents the architecture, the data schemas that define the contract between services, and the design decisions that shaped the platform. Every choice is grounded in the operational constraints of branch teller banking in the Tunisian regulatory context.

\section{System Overview}

The platform follows a \textbf{streaming-first} architecture: four microservices chained via message topics, supported by a hot/cold data layer and a batch pipeline for profile rebuilding and model retraining. Each transaction passes through four stages --- ingestion of the raw event, enrichment with behavioral context, scoring by two complementary models, and a decision layer merging the ML signal with deterministic regulatory rules. The output is a tiered alert with a human-readable explanation, consumed by a supervisor dashboard.

Three foundational choices define the architecture:

\begin{enumerate}
  \item \textbf{RedPanda over Apache Kafka.} A Kafka-compatible API without the JVM runtime or ZooKeeper dependency, reducing operational complexity at a throughput --- hundreds of events per minute from 150 branches --- that does not require Kafka's horizontal scaling. Critically, the architecture remains language-agnostic at topic boundaries: any service can be rewritten in a JVM language without changing the topic contracts.
  \item \textbf{Dual-actor profiling.} Client and processing employee are independently profiled, each with a hot profile in Redis and a cold profile in PostgreSQL. This is motivated by CTAF case studies in which collusion between a client and a specific teller was the key pattern.
  \item \textbf{SHAP co-located in the Scoring Service.} Explainability is computed at scoring time, not display time, and travels as part of the event payload. This reduces the dashboard to a single-topic consumer and guarantees that every alert carries its explanation regardless of how it is consumed downstream.
\end{enumerate}

\section{Streaming Topology}

Four RedPanda topics form the backbone, all partitioned by \code{client\_id}:

\begin{center}
\begin{tikzpicture}[node distance=0.55cm, every node/.style={font=\small}]
  \tikzstyle{topic}=[draw, rounded corners=2pt, minimum height=0.75cm, inner xsep=5pt]
  \node[topic] (a) {raw-events};
  \node[topic, right=of a] (b) {enriched-events};
  \node[topic, right=of b] (c) {scored-events};
  \node[topic, right=of c] (d) {decisions};
  \draw[-Stealth] (a) -- node[above,font=\tiny]{Enrich} (b);
  \draw[-Stealth] (b) -- node[above,font=\tiny]{Score} (c);
  \draw[-Stealth] (c) -- node[above,font=\tiny]{Decide} (d);
\end{tikzpicture}
\end{center}

Partitioning by \code{client\_id} is a deliberate constraint driven by the LSTM: it treats each client's history as a temporal sequence, so cross-operation ordering must be preserved per client. A \emph{retrait} followed by a \emph{versement} followed by a \emph{virement} must reach the Scoring Service in exactly that order, which partitioning by client guarantees.

For the same reason, a \textbf{single} \code{raw-events} topic carries all four operation types using a tagged-union envelope. Four separate topics --- one per operation --- were rejected: if \emph{retrait} and \emph{virement} events travel on separate topics there is no guarantee of relative ordering at the consumer.

A fifth topic, \code{dlq}, receives events that failed processing after retry exhaustion, preserving the original event alongside error metadata. Two sink consumers, on independent consumer groups, persist raw events and decisions to PostgreSQL, decoupling persistence from the transform-and-forward pipeline.

\section{Hot/Cold Data Layer}

\textbf{Redis} (hot) stores everything that must be available with sub-millisecond latency during enrichment: client and employee profiles, archetype baselines, per-client rolling event buffers for velocity features, and LSTM sequence buffers. \textbf{PostgreSQL} (cold) is the source of truth: all transactions, alert history, profile snapshots for audit and retraining, and master reference data. It backs the batch pipeline and the dashboard's queries.

The nightly batch pipeline reads PostgreSQL, recomputes aggregate statistics and behavioral distributions for both actors, writes snapshots back for audit, and refreshes the Redis hot profiles. The streaming services therefore never need to maintain complex windowed aggregations.

One subtlety: LSTM sequence buffers store \emph{unscaled} feature vectors, with the Scoring Service applying the current scaler at inference time. A weekly retrain produces a new scaler; storing pre-scaled vectors would invalidate every persisted sequence and force a full buffer rebuild on each retrain.

\section{Raw Event Schema}

The raw event uses a \textbf{tagged-union envelope}: 11 common fields shared by all operations plus an operation-specific \code{payload}. Two principles govern it. First, raw events contain \emph{only facts observable at transaction time} --- no aggregate statistics, no derived fields; statistics belong in profiles and derived features in the enriched event. Second, \emph{no lagging fields}: a cheque's bounce status is only known after clearing and belongs to a separate downstream event.

\begin{table}[htbp]
\centering
\small
\caption{Common envelope fields}
\begin{tabular}{@{}lll@{}}
\toprule
\textbf{Field} & \textbf{Type} & \textbf{Description} \\
\midrule
\code{event\_id} & UUID & Unique event identifier \\
\code{client\_id} & UUID & Client initiating the transaction \\
\code{account\_id} & UUID & Target account \\
\code{employee\_id} & UUID & Teller processing the transaction \\
\code{branch\_id} & String & Branch where the operation occurs \\
\code{timestamp} & Datetime & UTC timestamp \\
\code{amount} & Float & Transaction amount \\
\code{currency} & String & ISO 4217 code (typically TND) \\
\code{channel} & Enum & \code{guichet}, \code{GAB}, \dots \\
\code{operation\_type} & Enum & \code{retrait}, \code{versement}, \code{virement}, \code{cheque} \\
\code{payload} & JSON & Operation-specific data \\
\bottomrule
\end{tabular}
\end{table}

A \emph{virement} payload carries \code{beneficiary\_id} and \code{motif}; a \emph{cheque} payload carries \code{emitter\_id}, \code{cheque\_number}, and a \code{remise\_id} grouping cheques deposited together. Each cheque generates its own event rather than one event per \emph{remise}, preserving per-cheque granularity for scoring while keeping the grouping context.

\section{Dual-Actor Profiling}

The \textbf{client profile} comprises roughly 194 fields across six categories, summarized in Table~\ref{tab:profile}.

\begin{table}[htbp]
\centering
\small
\caption{Client profile categories}
\label{tab:profile}
\begin{tabular}{@{}p{3.2cm}p{1.4cm}p{8.5cm}@{}}
\toprule
\textbf{Category} & \textbf{Size} & \textbf{Content} \\
\midrule
Personal information & 7 & Archetype, maturity status, home branch, client type, account opening date \\
Transaction statistics & 113 & Count, sum, mean, std, min, max for each operation (4) $\times$ window (24h, 7d, 30d, 90d) \\
Behavioral patterns & 12 & Five frequency distributions (branch, hour, day-of-week, operation mix, employee) at 30d and 180d; known beneficiaries and emitters over 365d \\
Financial state & 12 & Balance statistics, inflow/outflow, net flow at 30d and 90d \\
Risk history & 49 & Alert counts per tier $\times$ window, each split into confirmed, rejected, pending \\
Sequence buffer & 50 events & Bounded Redis list matching the LSTM training sequence length \\
\bottomrule
\end{tabular}
\end{table}

Jensen--Shannon divergence between the 30-day and 180-day distributions captures behavioral drift in a single scalar per dimension. The risk history split by supervisor verdict enables both recidivist escalation and false-positive dampening in the Decision Service.

The \textbf{employee profile} mirrors this structure with two differences. Financial state is excluded --- monitoring an employee's personal finances is out of scope and ethically fraught. A \emph{peer comparison} category is added: five z-scores at 30d and 90d covering volume, error rate, near-threshold ratio (fraction of transactions in the 8{,}000--9{,}999~DT band), client concentration, and client-location mismatch.

A normalization challenge follows from the branch structure: with only 2 tellers per branch, per-branch z-scores are statistically meaningless. The solution is \textbf{split normalization}. Volume metrics are first normalized per-branch (employee volume divided by branch average), then z-scored bank-wide --- preventing high-volume branches from appearing as outliers while still detecting employees handling disproportionate volume within their branch. Ratio metrics are z-scored directly bank-wide, since ratios are already volume-independent.

\section{Enriched Event Design}

The enriched event is the feature vector consumed by the Scoring Service. Its redesign produced the single most impactful improvement in detection performance.

\subsection{The Profile Dominance Problem}

The initial schema contained 74 features, many of them raw profile field dumps --- weekly snapshot means, standard deviations, and full distribution vectors repeated identically across every event from the same client within the same week. The per-event signal (how unusual \emph{this} transaction is) was drowned out by profile noise (what the client's long-term averages are). Both models failed to separate normal from anomalous events, producing AUC near random ($\sim$0.57).

\subsection{Contrast Feature Redesign}

The solution was to replace profile dumps with \textbf{contrast features}: each feature measures the deviation of the current event against the client's established baseline. The count dropped from 74 to 30.

\begin{table}[htbp]
\centering
\small
\caption{Enriched event feature categories (30 features; full list in Appendix~\ref{app:features})}
\begin{tabular}{@{}p{3.6cm}c p{8.2cm}@{}}
\toprule
\textbf{Category} & \textbf{Count} & \textbf{Representative features} \\
\midrule
Raw event & 7 & amount, hour, operation one-hot, \code{is\_home\_branch} \\
Amount contrast & 6 & \code{z\_amount}, ratio to mean, threshold and round-amount flags \\
Operation contrast & 1 & \code{op\_type\_probability} from the 30d distribution \\
Timing contrast & 2 & \code{is\_outside\_hours}, day distance from preferred day \\
Velocity & 8 & 24h/7d counts, cumulative amount, duplicates, near-threshold count, employee counters \\
Context & 6 & archetype one-hot, \code{account\_age\_days} \\
\bottomrule
\end{tabular}
\end{table}

The key invariant: the \code{\_compute\_features()} function producing these 30 features is a \textbf{pure function} --- identical inputs always produce identical outputs, with no I/O and no side effects. Both the streaming pipeline and the batch training pipeline call this same function, guaranteeing zero feature drift between training and inference. This is the single most important technical invariant in the system.

\section{Scoring Architecture}

\subsection{Dual-Model Rationale}

The \textbf{Autoencoder}, trained on normal transactions only, flags events with high reconstruction error. It excels at \emph{point anomalies} --- a 15{,}000~DT \emph{retrait} from a salaried client averaging 300~DT --- and works from the very first event, with no history required.

The \textbf{LSTM} predicts what a client's next event should look like from their recent sequence. It excels at \emph{sequential anomalies} --- a gradual increase in \emph{versement} frequency where each individual event looks normal but the trajectory signals structuring. It requires sufficient history to be meaningful.

\subsection{Score Fusion}

The two raw scores live on incompatible scales (reconstruction MSE versus prediction MSE). Three mechanisms reconcile them:

\begin{enumerate}
  \item \textbf{Percentile normalization.} Each raw score is mapped to a percentile rank against a reference distribution computed from normal training data, placing both on a common $[0,1]$ scale where 0.95 means \emph{``higher than 95\% of normal events.''}
  \item \textbf{Maturity-adaptive weighting.} A sigmoid of \code{days\_since\_account\_opening} (midpoint 90 days) controls the balance. This was chosen over event count, which penalizes low-frequency archetypes: a retiree transacting twice a month is not ``new'' but has a low event count.
  \item \textbf{Weighted combination:} $S_{\text{fused}} = w_{\text{lstm}} \cdot S_{\text{lstm}} + (1 - w_{\text{lstm}}) \cdot S_{\text{ae}}$.
\end{enumerate}

\subsection{Cold-Start Scoring}

Clients with fewer than 5 events in their sequence buffer fall outside the training distribution entirely, since the training pipeline applies a 21-day warm-up filter. A three-tier progression handles this:

\begin{table}[htbp]
\centering
\small
\caption{Cold-start scoring progression}
\begin{tabular}{@{}lll@{}}
\toprule
\textbf{Buffer length} & \textbf{Scoring behavior} & \textbf{Tier cap} \\
\midrule
0--4 events & Neutral score (0.5), \code{cold\_start=True} & REVIEW max \\
5--9 events & AE only ($w_{\text{lstm}} = 0$) & Full range \\
10+ events & Full AE + LSTM fusion & Full range \\
\bottomrule
\end{tabular}
\end{table}

Cold-start events still pass through all regulatory rules --- a 15{,}000~DT transaction must trigger the BCT reporting flag regardless of buffer length --- but the ML-derived tier is capped at REVIEW to prevent score saturation from generating false BLOCKs.

\subsection{Conditional SHAP}

SHAP is expensive ($\sim$100\,ms per event via KernelExplainer) and roughly 95\% of events score below the REVIEW threshold, so computation is gated behind the threshold check: below REVIEW the event takes the fast path ($\sim$5\,ms). A deterministic template dictionary maps feature names to parameterized sentences, producing consistent and auditable explanations --- \emph{``This versement is suspect because the amount (8{,}450~DT) is 4.2$\times$ the client's average, and it is the 3rd deposit in 48 hours.''} An LLM-based generator was rejected on bank security grounds and for non-determinism.

\section{Decision Service}

The Decision Service is a stateless three-layer engine merging the ML signal with deterministic rules and historical context.

\textbf{Layer 1 --- score thresholds.} Three thresholds calibrated on the validation set define the tiers: REVIEW at $T_1 = 0.95$ (review within 24 hours), ALERT at $T_2 = 0.99$ (immediate attention), BLOCK at $T_3 = 0.999$ (transaction held, requires override). Events below $T_1$ receive INFO: logged for audit, no action required.

\textbf{Layer 2 --- regulatory escalation.} Eight rules can escalate the tier from Layer 1 but \emph{never de-escalate} it. This is a hard invariant: a regulatory flag cannot reduce alert severity, and all flags are logged for audit whether or not they change the tier.

\begin{table}[htbp]
\centering
\small
\caption{Regulatory escalation rules}
\begin{tabular}{@{}rp{3.3cm}p{6.4cm}l@{}}
\toprule
\# & \textbf{Rule} & \textbf{Condition} & \textbf{Escalation} \\
\midrule
1 & BCT reporting & Amount $\geq$ 10{,}000~DT and $z_{\text{amount}} \geq 2$ & ALERT \\
2 & Smurfing band & Amount in 8{,}000--9{,}999~DT and $\geq$2 near-threshold txns in 7d & ALERT \\
3 & Cheque ceiling & Amount in 25{,}000--30{,}000~DT (Law n\textsuperscript{o}41-2024) & REVIEW \\
4 & Passeur de fonds & Occasional client, $\geq$5 versements in 24h & ALERT \\
5 & New beneficiary & Virement $\geq$ 5{,}000~DT to first-time beneficiary & REVIEW \\
6 & Duplicate & Matching operation within time window & REVIEW \\
7 & Round amount & Amount $\geq$ 5{,}000~DT, round, high 24h cumulative & REVIEW \\
8 & Frequency burst & $\geq$ 5$\times$ expected daily rate & REVIEW \\
\bottomrule
\end{tabular}
\end{table}

Rule~1's $z_{\text{amount}} \geq 2$ guard matters: a 15{,}000~DT \emph{virement} from a big-business client averaging 20{,}000~DT is routine and should not flood the REVIEW queue. Only transactions that are \emph{both} above the regulatory threshold \emph{and} unusual for the specific client escalate.

\textbf{Layer 3 --- risk history adjustment.} Clients with $\geq$3 confirmed alerts in 30 days receive a tier boost (a borderline event from a repeat offender is treated more seriously), while new accounts receive a more conservative posture reflecting the higher risk of unestablished baselines.

Explanation assembly depends on what triggered the alert: ML-only (SHAP contributions plus LSTM expected-vs-actual), rule-only (rule description with the specific values), or both (a combined explanation).

\section{Software Architecture}

A mid-project requirement --- a synchronous REST endpoint for the dashboard's simulation mode --- exposed a coupling problem: all business logic lived inside Kafka consumer/producer wrappers, and a synchronous \code{/score} request cannot route through an asynchronous pipeline without unacceptable complexity. Request--reply over Kafka was rejected on latency; duplicating the logic outside Kafka was rejected on maintainability. The chosen option was an in-process orchestration layer, implemented through a \textbf{core extraction pattern}:

\begin{itemize}
  \item \textbf{Core module} (\code{EnrichmentCore}, \code{ScoringCore}, \code{DecisionCore}): receives external dependencies --- Redis client, model objects, metrics collector --- via dependency injection, exposes a single public method (fa\c{c}ade pattern), contains all business logic, and has \emph{no} Kafka dependency.
  \item \textbf{Streaming wrapper}: thin consumer/producer plumbing that receives a core instance via DI; its poll loop calls the core's public method and produces the result. Contains no business logic.
\end{itemize}

The FastAPI simulation API imports the three cores directly, wires them with shared Redis and model dependencies at startup via a lifespan context manager, and exposes a synchronous \code{/score} endpoint calling \emph{enrich $\to$ score $\to$ decide} in-process. A \code{dry\_run} flag on enrichment prevents simulation requests from polluting the streaming pipeline's buffer state. The same pattern lets the batch pipeline call \code{\_compute\_features()} directly, guaranteeing feature parity across three execution contexts: streaming, batch, and simulation.
